\section{Separate the initial byte into two fields}
An item begins with an initial byte. Its upper three bits select a major type; its
lower five bits supply additional information. Read the byte unsigned before shifting
or masking it. For the worked input's first byte, hexadecimal \texttt{b8}, the upper
field is five and the lower field is twenty-four. The latter tells the reader to obtain
the argument from a following byte, which is \texttt{04} in this input.

\begin{figure}[htbp]
\centering
\begin{tikzpicture}[x=0.85cm,y=0.6cm,alt={One initial byte has eight bits. From left to right they are bits seven through zero. The first three bits form the major type; the last five form additional information. For hexadecimal b8 these groups are binary 101 and 11000.}]
\foreach \i in {0,...,7}{
 \pgfmathtruncatemacro{\bit}{7-\i}
 \ifnum\i<3 \draw[fill=black!12] (\i,0) rectangle +(1,1);\else \draw (\i,0) rectangle +(1,1);\fi
 \node at (\i+0.5,0.5) {\bit};
}
\node[below] at (1.5,-0.15) {major type};
\node[below] at (5.5,-0.15) {additional information};
\end{tikzpicture}
\caption{The two fields of an initial byte. This is a native format diagram, not an IDE capture.}
\end{figure}

Together, the initial byte and any bytes carrying its argument form the item's head.
The head does not include a byte string's contents or an array's children. Here the
head occupies two bytes and describes a map with four pairs. Returning four as the
head's byte length would advance the caller into the first optional value.
